\chapter{How the Code Works}

\section{The map}

Twenty-odd Python files, but only six carry the science. Everything else is
orchestration or diagnosis.

\begin{center}
\begin{tabular}{@{} L{4.4cm} L{9.4cm} @{}}
\toprule
\textbf{File} & \textbf{Responsibility} \\
\midrule
\multicolumn{2}{@{}l}{\textbf{Core --- the science}}\\
\texttt{arm.xml} & the physical model: bodies, joints, muscles, geometry \\
\texttt{config.py} & every constant in one place; single source of truth \\
\texttt{rl/environment.py} & the task: observation, reward, episode rules \\
\texttt{rl/train.py} & the training loop, checkpointing, learning curve \\
\texttt{rl/evaluate.py} & rollout and metrics --- produces every reported number \\
\texttt{force\_comparison.py} & \textbf{the central analysis}: DRL vs classical \\
\midrule
\multicolumn{2}{@{}l}{\textbf{Supporting analyses}}\\
\texttt{run\_conventional.py} & a classical controller solving the same task \\
\texttt{run\_synergies.py} & muscle-synergy decomposition \\
\texttt{analysis/synergies.py} & the NMF machinery \\
\texttt{bench\_loadsharing.py} & like-for-like latency measurement \\
\midrule
\multicolumn{2}{@{}l}{\textbf{Verification --- written to catch specific failures}}\\
\texttt{verify\_replay\_ratio.py} & measures gradient updates per environment step \\
\texttt{verify\_divergence.py} & audits simulator warnings against reported blow-ups \\
\texttt{test\_reward\_v3.py} & 14 checks on the reward's properties \\
\texttt{oracle\_feasibility.py} & can the muscles hold the arm at all? \\
\midrule
\multicolumn{2}{@{}l}{\textbf{Orchestration}}\\
\texttt{run\_parallel.py} & runs jobs across CPU and GPU concurrently \\
\texttt{reevaluate.py} & re-runs all evaluations after the evaluation-code fixes \\
\bottomrule
\end{tabular}
\end{center}

\section{How the parts talk to each other}

\begin{center}
\begin{tabular}{@{} L{15cm} @{}}
\toprule
\ttfamily\small
config.py \ \ (constants: weights, budgets, seeds, paths)\\
\ttfamily\small
\ \ \ $\downarrow$ read by everything\\[4pt]
\ttfamily\small
arm.xml $\rightarrow$ environment.py $\rightarrow$ [obs 38] $\rightarrow$ algorithms.py $\rightarrow$ SB3 model\\
\ttfamily\small
\ \ \ \ \ \ \ \ \ \ \ \ $\uparrow$ [action 9] \ \ \ \ \ \ \ \ \ \ \ \ \ \ \ \ \ \ \ \ \ \ \ \ \ \ \ \ $\downarrow$\\
\ttfamily\small
\ \ \ \ \ \ \ \ \ \ \ \ +---------------- train.py (loop, checkpoints) -----+\\[4pt]
\ttfamily\small
train.py $\rightarrow$ runs/<name>/\{model.zip, vecnormalize.pkl, train\_meta.json\}\\
\ttfamily\small
\ \ \ \ \ \ \ \ \ \ \ \ \ \ \ \ \ \ \ \ $\downarrow$\\
\ttfamily\small
evaluate.py $\rightarrow$ metrics \ \ \ force\_comparison.py $\rightarrow$ effort ratio, per-muscle\\
\ttfamily\small
\ \ \ \ \ \ \ \ \ \ \ \ \ \ \ \ \ \ \ \ $\downarrow$ \ \ \ \ \ \ \ \ \ \ \ \ \ \ \ \ \ \ \ \ \ \ \ \ \ \ \ \ \ $\downarrow$\\
\ttfamily\small
\ \ \ \ \ \ \ \ \ \ \ \ \ \ \ \ \ \ runs/*.json $\rightarrow$ build\_figures*.py $\rightarrow$ the report\\
\bottomrule
\end{tabular}
\end{center}

\begin{why}
Two design rules hold throughout, and both were adopted after being violated.

\textbf{One source of truth.} Every constant lives in \texttt{config.py}. Nothing
is written twice. When the effort weight changed, it changed in one place and
every downstream script saw it.

\textbf{Numbers reach the report only through files.} No figure or table is typed
by hand. Scripts write JSON; \texttt{build\_figures*.py} reads that JSON. This is
a direct response to discovering that the previous version of this thesis
contained numbers generated by \texttt{numpy.random.normal} and never replaced.
\end{why}

\section{\texttt{rl/environment.py} --- the task}

This is where the problem is actually defined. Everything else is machinery.

\subsection{The observation: what the agent sees}
\label{sec:obs}

38 numbers, assembled fresh each step:

\begin{center}
\begin{tabular}{@{} C{1.8cm} C{1.2cm} L{5.0cm} L{6.2cm} @{}}
\toprule
\textbf{Indices} & \textbf{Count} & \textbf{Content} & \textbf{Why the agent needs it} \\
\midrule
0--3   & 4 & joint angles $q$ (rad) & where the arm is \\
4--7   & 4 & joint velocities $\dot q$ (rad/s) & whether it is moving, and how fast --- required to decelerate \\
8--16  & 9 & muscle lengths & determines available force via force--length \\
17--25 & 9 & muscle forces (N) & what each muscle is currently doing \\
26--34 & 9 & muscle activations & the internal state lagging the command \\
35--37 & 3 & target minus hand position & where to go, as a vector \\
\bottomrule
\end{tabular}
\end{center}

\begin{why}
Two choices here are deliberate.

The target is given as a \emph{difference} (target minus hand), not as an
absolute coordinate. The agent therefore learns ``move in this direction'' rather
than ``go to this place'', which generalises across the workspace instead of
memorising locations.

Joint velocities are included because without them the task is not solvable. A
controller that cannot perceive that it is moving cannot decide to stop --- and
stopping on target is precisely what this task requires.
\end{why}

\subsection{The reward}
\label{sec:reward}

\begin{multline*}
r \;=\; \underbrace{w_4}_{\text{survival}}
\;-\;\underbrace{w_1\big(\|e\|^2 + 0.5\|e\|\big)}_{\text{distance cost}}
\;-\;\underbrace{w_2\, \overline{(f_i/F_{\max,i})^2}}_{\text{effort}}
\;-\;\underbrace{w_3\, \overline{(\Delta a_i)^2}}_{\text{smoothness}}\\[4pt]
\;+\;\underbrace{w_7 e^{-\|e\|/s}}_{\text{precision bonus}}
\;+\;\underbrace{w_5 \,[\text{on target}]}_{\text{success}}
\;-\;\underbrace{w_6 \,[\text{blow-up}]}_{\text{penalty}}
\end{multline*}

Term by term:

\begin{longtable}{@{} L{2.6cm} C{1.4cm} L{10.4cm} @{}}
\toprule
\textbf{Term} & \textbf{Weight} & \textbf{What it does and why it exists} \\
\midrule
\endhead
Survival & $w_4 = 4.5$ & A constant paid every step. Makes the per-step reward strictly positive, so ending an episode early forfeits value. Added specifically to kill the self-destruction exploit. A constant cannot reorder trajectories of equal length, so it does not disturb the other terms' tuning. \\
Distance & $w_1 = 1.0$ & Quadratic \emph{and} linear in error. Quadratic alone has no gradient near the target ($e^2 = 0.01$ at 10\,cm), so policies hovered; the linear term keeps a pull at all ranges. \\
Effort & $w_2$ & Mean squared force as a fraction of each muscle's maximum. This is the classical criterion in form. \textbf{Its weight is the subject of a whole experiment} (\S\ref{sec:sweepresult}). \\
Smoothness & $w_3 = 0.5$ & Penalises abrupt changes in command, discouraging jitter. Mean, not sum, so it is bounded by 1 like the effort term. \\
Precision bonus & $w_7 = 3.0$ & $e^{-\|e\|/s}$ with $s = 0.15$\,m. Negligible far away, steep close in --- a dense, shaped version of the sparse success bonus, added because the sparse one was never once collected. \\
Success & $w_5 = 1.0$ & Paid \emph{every step} the hand is on target and settled, not once on arrival. This makes the task reach-\emph{and}-hold. \\
Blow-up & $w_6 = 10.0$ & One-off penalty for numerical divergence. \\
\bottomrule
\end{longtable}

\begin{trap}
The precision bonus was originally $s = 0.05$\,m. At the \SI{0.23}{\meter} where
the policy actually sat, $e^{-0.23/0.05} = 0.010$ --- the bonus contributed
$0.03$ of a reward of about 4, i.e.\ nothing. It only started paying once the hand
was already within \SI{15}{\centi\meter}, a region the policy could not reach. It
was a prize behind a wall. Setting $s = 0.15$ matched the scale to where the
policy actually lived.
\end{trap}

\subsection{Episode rules}

\begin{itemize}[leftmargin=1.4em, itemsep=2pt]
\item \textbf{Truncation} at 100 agent steps (2 seconds).
\item \textbf{Termination} only on numerical divergence.
\item \textbf{Success does not end the episode} --- arriving early means
collecting the bonus for staying, which is the point.
\item \textbf{Targets} are generated by picking random joint angles and computing
where the hand would be. Every target is therefore reachable by construction ---
an earlier version sampled from a box, much of which lay outside the workspace.
\end{itemize}

\subsection{Action repeat (frame skip)}

The policy acts at \SI{50}{\hertz} while physics runs at \SI{500}{\hertz}: each
command is held for 10 physics steps.

\begin{why}
This looks like a detail and is worth an experiment's difference. At one action
per physics step, a 2-second episode is 1000 decisions, so a million-step budget
buys only 1000 reaching \emph{attempts} --- and the policy regressed toward the
middle of the workspace from sheer lack of experience.

At \SI{50}{\hertz} the same budget buys about 10{,}000 attempts. And it is nearly
free: a physics step costs \SI{6.8}{\micro\second} while a gradient update costs
about \SI{14}{\milli\second}, so wall-clock is set by learning, not simulation.
Ten times the physics adds roughly a minute to a 75-minute run.
\end{why}

\section{\texttt{rl/train.py} --- the training loop}

Wraps Stable-Baselines3 with three additions.

\textbf{Learning curve.} Every 25{,}000 or 50{,}000 steps, run five deterministic
episodes and record error, success, blow-up rate and episode length. It is a
monitoring signal only; no reported number comes from it, because five episodes
is far too noisy.

\textbf{Checkpointing.} Every 50{,}000 steps, save model weights, normalisation
statistics, the replay buffer and progress. Written to a temporary file and
atomically renamed, so an interrupted save cannot corrupt the previous one.

\begin{trap}
This exists because a forced operating-system restart destroyed a 75-minute run
five minutes from completion. The original code saved only after training
finished, so nothing was recoverable. Atomic renaming matters because a crash
\emph{during} a save would otherwise leave an unloadable half-written file where
a good checkpoint used to be.
\end{trap}

\textbf{Best-model tracking.} Saves whichever checkpoint scored best on a
validation stream seeded differently from the reporting stream, so the reported
number is not the one selection was performed on. (In practice all reported
results use the \emph{final} model, so this introduces no selection bias.)

\section{\texttt{rl/evaluate.py} --- where the numbers come from}

Loads a trained policy, runs 50 deterministic episodes with domain randomisation
off, and computes every metric.

\begin{center}
\begin{tabular}{@{} L{3.4cm} L{10.6cm} @{}}
\toprule
\textbf{Metric} & \textbf{Definition and purpose} \\
\midrule
final error & hand-to-target distance at the end. The headline accuracy number. \\
closest approach & smallest distance at any instant. Separates ``could it get there'' from ``could it stop there''. \\
drift & final minus closest. How far it wandered back out. \\
reach 5/10\,cm & fraction of episodes \emph{ending} within that distance \\
touch 2/5\,cm & fraction reaching within that distance at any moment \\
energy & $\sum_i F_i^2$, a metabolic proxy \\
blow-up rate & fraction of episodes ending in divergence \\
episode length & mean steps survived; \textbf{exists to catch an agent that quits} \\
success rate & within 2\,cm \emph{and} nearly stationary at the end \\
\bottomrule
\end{tabular}
\end{center}

\begin{why}
The last two exist because of the self-destruction exploit. When it was
occurring, the recorded error was still a plausible, slowly improving number, and
nothing logged how \emph{long} episodes lasted --- so an agent quitting after 7
of 100 steps looked identical to one trying hard and plateauing. Episode length
and blow-up rate make that failure impossible to hide.

\textbf{Success rate is deliberately not a headline metric.} The criterion is one
that a controller \emph{with privileged information} satisfies only about 7\,\%
of the time. A measure that a near-best-possible controller fails 93\,\% of the
time cannot discriminate between learned policies. Across five identically
configured runs it read 0\,\%, 4\,\%, 6\,\%, 8\,\% and 20\,\% --- a twentyfold
spread that identifies it as noise.
\end{why}

\begin{trap}
Four defects were later found in this file, none producing arithmetically wrong
numbers but two changing \emph{which episodes} the numbers describe:
the evaluation seed was dead code; the environment was reset twice per episode,
so evaluation scored on every \emph{other} target; the device setting was
ignored; and a diverged episode could be scored a success. All are documented in
\S\ref{sec:evaldefects2} with what they changed.
\end{trap}

\section{\texttt{force\_comparison.py} --- the central analysis}

This file answers the thesis question. Everything else supports it.

At every timestep of a policy rollout it solves the classical problem on the
movement the policy just produced:
\[
\min_{f}\ \sum_i \left(\frac{f_i}{F_{\max,i}}\right)^2 \quad\text{s.t.}\quad R^\top f = \tau,\quad 0 \le f_i \le F_{\max,i}
\]
and compares $f_{\text{classical}}$ against the $f_{\text{policy}}$ MuJoCo
actually produced.

\begin{why}
\textbf{Where does $\tau$ come from?} It is taken as \texttt{qfrc\_actuator} ---
the generalised force the muscles actually produced at that instant.

Two reasons. It is \emph{exact}, avoiding the numerical differentiation that
inverse dynamics would introduce. And it guarantees a solution exists, because
the policy's own force vector already satisfies the constraint --- it is a
witness. The optimiser can therefore never be asked something impossible.

The consequence is that both methods are posed an \emph{identical} problem on an
\emph{identical} trajectory, so any difference is purely load-sharing strategy.
\end{why}

\subsection{The validity check, and why it is not optional}

The whole comparison depends on $R^\top f = \tau$ holding for the policy's own
forces. If the moment-arm extraction or the sign convention were wrong, every
number would be meaningless while still looking reasonable.

So it is checked and printed every run. Measured: mean residual
\SI{1.46e-15}{\newton\meter} --- machine precision --- and 0 optimiser failures
in 2000 timesteps.

\subsection{The metrics it produces}

\begin{center}
\begin{tabular}{@{} L{3.6cm} L{10.4cm} @{}}
\toprule
\textbf{Metric} & \textbf{Meaning} \\
\midrule
effort ratio & total policy effort $\div$ total classical effort. \textbf{The headline.} 1.0 = the policy found the minimum-effort solution. \\
Pearson $r$ & correlation between the two force sets. Mean-centred, so no artificial floor. \\
pattern cosine & similarity of direction of the 9-vector, ignoring magnitude \\
null baseline & the same cosine with muscle identities shuffled \\
per-muscle RMSE & where the disagreement lives, anatomically \\
\bottomrule
\end{tabular}
\end{center}

\begin{trap}
\textbf{Cosine similarity must never be quoted alone.} Muscle forces are all
non-negative, so two unrelated force vectors already agree substantially by
construction. The null --- shuffling which muscle each force belongs to, keeping
magnitudes --- scores $0.373$. The observed $0.869$ is therefore an excess of
$+0.496$, not a near-perfect $0.87$.

Worse, the null's 95th percentile for \emph{individual} timesteps is $0.867$, so
no single timestep's cosine means anything; only the aggregate does. The null is
now computed automatically and written to the output file, so the bare number
cannot reappear.
\end{trap}

\section{\texttt{run\_conventional.py} --- the classical baseline}
\label{sec:conventional}

Static optimisation only ever answers the load-sharing sub-question --- it is
handed a movement. To ask whether classical control could do the \emph{whole
job}, a conventional controller was built:

\begin{enumerate}[leftmargin=1.5em, itemsep=2pt]
\item Compute a desired hand force: $f_{\text{des}} = K_p(x^\star - x) - K_d\dot{x}$.
\item Map it to joint torques through the Jacobian transpose.
\item \textbf{Add gravity, Coriolis and centrifugal compensation.}
\item Probe each muscle's achievable force range at the current state.
\item Solve minimum-effort load sharing over that range.
\item Convert the resulting forces back to activations.
\end{enumerate}

\begin{trap}
Steps 3 and 4 were both absent initially, and each alone made the baseline a
strawman.

Without gravity compensation the arm droops until position error alone generates
enough torque to hold it --- a large steady-state offset. Measured: it reached
only \SI{0.46}{\meter}, \emph{worse than a random policy}.

Without the achievable-range probe, the code assumed force is proportional to
activation. It is not: force is $\text{gain}(l,v)\cdot a + \text{bias}(l,v)$,
state-dependent. Converting a desired \emph{force} to an \emph{activation}
requires inverting the muscle model --- exactly the operation this thesis
proposes to replace. The baseline was skipping the step that makes it
conventional.

Publishing ``DRL beats classical control'' on the uncorrected version would have
been a strawman comparison. Both corrections improve the \emph{classical} side.
\end{trap}
